\documentclass[twocolumn,superscriptaddress,aps,prl,floatfix]{revtex4-2}

% ============================================================
% DiffCMB — Draft manuscript
% "A full-sky differentiable joint sampler over the CMB lensing potential"
%
% DRAFTING RULES (ROADMAP.md, 2026-08-06):
%   - NEVER lead with the differentiable-SHT machinery. Flinch (arXiv:2510.26691)
%     made curved-sky differentiable inference table stakes. Every abstract, intro
%     sentence and talk title leads with the JOINT (C_ell, C_L^phiphi) POSTERIOR.
%   - Differentiable infrastructure goes in Methods as "enabling machinery,
%     consistent with Flinch", never in the novelty sentence.
%   - Every "first" carries scope qualifiers and nearest-prior-work citations.
%   - Re-check the named-author arXiv scan (Millea, Seljak, Bayer, Loureiro)
%     before every submission milestone.
%
% SECTION STATUS:
%   [DONE]        = content fixed, does not depend on pending results
%   [PLACEHOLDER] = leave stub; will fill once chains are done
%   [TODO]        = not started; needed before submission
% ============================================================

\usepackage{amsmath,amssymb,bm}
\usepackage{graphicx}
\usepackage{hyperref}
\usepackage{xcolor}
\usepackage{natbib}

% Convenience macros
\newcommand{\alm}{a_{\ell m}}
\newcommand{\cl}{C_\ell}
\newcommand{\clpp}{C_L^{\phi\phi}}
\newcommand{\phiphi}{\phi\phi}
\newcommand{\lmax}{\ell_{\max}}
\newcommand{\todo}[1]{\textcolor{red}{\textbf{[TODO: #1]}}}
\newcommand{\placeholder}[1]{\textcolor{blue}{\textbf{[PLACEHOLDER: #1]}}}

% ============================================================
\begin{document}

\title{A full-sky differentiable joint Gibbs sampler over
       the CMB temperature field, power spectrum, and lensing potential}

\author{\placeholder{Oscar Hickman}}
\affiliation{\placeholder{Institute for Computational Cosmology, Durham University}}

% Additional authors to be confirmed
\date{\today}

% ============================================================
\begin{abstract}
% RULE: lead with the joint posterior object, not the machinery.
% Draft language from ROADMAP.md 2026-08-06 re-pitch:
We present the first curved-sky, full-sky posterior over the joint
$(\alm, \cl, \phi, \clpp)$ — a four-block Gibbs sampler operating directly on
HEALPix maps whose blocks alternate exact inverse-Gamma draws for the two power
spectra and Hamiltonian Monte Carlo / microcanonical Langevin steps for the two
field amplitudes.
The sampler's defining output is the joint $(\cl^{TT}, \clpp)$ posterior with
correlations propagated rather than assumed, an object no existing method
produces: quadratic estimators and iterative MAP reconstruction (e.g.\
\citealt{Carron:2017, MUSE}) return point estimates or marginal constraints;
Commander-class samplers \citep{Commander} and their differentiable descendants
\citep{Almanac,Flinch} infer maps and power spectra on the full sky but carry no
lensing-potential block.
We validate the sampler via simulation-based calibration (SBC) on lensed
temperature maps at $\lmax = \placeholder{128}$, demonstrating frequentist
coverage of the joint posterior.
As an exact reference standard, the sampler's posterior is the natural
calibration target for fast, marginal, or learned CMB lensing methods.
\todo{Add key numbers once chains are done: coverage pass, joint-correlation amplitude.}
\end{abstract}

\maketitle

% ============================================================
\section{Introduction}
\label{sec:intro}
% STATUS: [DONE — draft from ROADMAP.md re-pitch; sharpen at writing time]

% --- Opening: what the competing landscape is ---
Curved-sky differentiable field-level CMB inference is now established.
\citet{Flinch} demonstrate gradient propagation from masked full-sky temperature
maps to cosmological parameters, recovering maps \textit{and} angular power
spectra;
\citet{Almanac} sample all-sky CMB maps jointly with their power spectra via
Hamiltonian Monte Carlo.
Both analyses are \textit{lensing-blind}: they infer $(\alm, \cl)$ and, in
Flinch's case, cosmological parameters, but carry no lensing-potential $\phi$ and
no $\clpp$ block.

Conversely, the methods that do treat lensing on the curved sky — quadratic
estimators, iterative maximum \textit{a posteriori} (MAP) reconstruction
\citep{Carron:2017}, and MUSE \citep{MUSE} — return point estimates or marginal
constraints, not a joint posterior.
The flat-sky code CMBLensing.jl \citep{CMBLensing} samples jointly over
$(\alm, \phi)$ on the flat sky but does not cover the full sky.

This paper closes the remaining cell: a full-sky, HEALPix, differentiable joint
Gibbs sampler over $(\alm, \cl, \phi, \clpp)$, validated by
simulation-based calibration.
Its primary output is the joint $(\cl^{TT}, \clpp)$ posterior and its
correlations — an object no existing method produces — which functions as a
\textit{reference standard}: the exact posterior that fast, marginal, or learned
methods \citep{Doeser:2024} get validated against.

% --- What we offer vs each competing paradigm ---
\subsection*{Position vs existing methods}

\textbf{vs Commander / Almanac / Flinch} (exact field samplers, lensing-blind):
our sampler adds the $(\phi, \clpp)$ blocks.
The differentiable spherical-harmonic infrastructure we use is consistent with
\citet{Flinch} and treated here as enabling machinery, not as a contribution.

\textbf{vs QE / iterative MAP / MUSE} (lensing-sensitive, marginal):
our sampler provides a joint posterior rather than a point estimate or a
marginalised constraint; it is not competitive in speed but provides
the calibration target these methods lack.

\textbf{vs diffusion / score-based generative methods} \citep{Drouin:2023}:
diffusion samplers produce approximate uncorrelated samples quickly but discard
the two distinguishing properties we build: a differentiable forward model and
a statistically exact sampler.
An exact sampler is what learned posteriors get validated against
\citep{Doeser:2024}.
The comparison is not speed but correctness.

\textbf{Honest scope statement.}
We demonstrate coverage at $\lmax \approx \placeholder{128}$.
Scaling to $\lmax \sim 1000$ and polarisation (the science-driver configuration)
is left to follow-up work.


% ============================================================
\section{Method}
\label{sec:method}
% STATUS: [DONE for block structure; sub-sections to be fleshed out]

We describe the four-block Gibbs sampler.
Let $\mathbf{d}$ denote the observed (lensed, beam-convolved, noise-added)
temperature map.
The joint target is
\begin{equation}
    p(\alm, \cl, \phi, \clpp \mid \mathbf{d})
    \propto p(\mathbf{d} \mid \alm, \phi)\,
            p(\alm \mid \cl)\,p(\cl)\,p(\phi \mid \clpp)\,p(\clpp),
\end{equation}
where $p(\mathbf{d} \mid \alm, \phi)$ is the lensed CMB likelihood evaluated via
a differentiable full-sky lensing operator.

% --- Block 1 ---
\subsection{Block 1: $\cl \mid \alm$ (exact inverse-Gamma draw)}
\label{sec:block1}

Given the current unlensed alm amplitudes, the conditional $p(\cl \mid \alm)$
is exactly inverse-Gamma (conjugate to the Gaussian alm prior) and is sampled
exactly by a one-line draw.
This is the standard Gibbs move used in Commander \citep{Commander} and
Almanac \citep{Almanac}.

% --- Block 2 ---
\subsection{Block 2: $\alm \mid \cl, \phi$ (HMC)}
\label{sec:block2}

We sample the unlensed alm given $\cl$ and $\phi$ via Hamiltonian Monte Carlo,
using the matrix-free ducc0 spherical-harmonic transform \citep{Reinecke:2023}
wrapped in a \texttt{tf.custom\_gradient} to propagate autodiff gradients through
the lensing operator.
This is enabling machinery consistent with \citet{Flinch}; the differentiable
SHT is not presented as a contribution.
The mass matrix is the posterior precision $1/\cl$; whitened coordinates are
used throughout.

% --- Block 3 ---
\subsection{Block 3: $\phi \mid \alm, \cl$ (MCLMC)}
\label{sec:block3}

We sample the lensing potential $\phi$ given the current alm and $\cl$ via
microcanonical Langevin Monte Carlo (MCLMC) \citep{Robnik:2022}, an unadjusted
isokinetic-leapfrog integrator.
MCLMC is chosen over standard HMC based on the evidence of
\citet{Bayer:2023,Flinch}, who report one-to-three-order-of-magnitude efficiency
gains over HMC at exactly this dimensionality regime on curved-sky CMB fields.
The MCLMC integrator is hand-implemented in TensorFlow, reusing the same
lensed log-probability target and its autodiff gradient unchanged.

% --- Block 4 ---
\subsection{Block 4: $\clpp \mid \phi$ (exact inverse-Gamma draw)}
\label{sec:block4}

Given the current $\phi$ amplitudes, the conditional $p(\clpp \mid \phi)$ is
exactly inverse-Gamma by the same conjugacy argument as Block 1,
\begin{equation}
  \clpp \mid \phi \;\sim\; \mathrm{InvGamma}\!\left(L - \tfrac12,\; S_L/2\right),
  \qquad S_L = \sum_{m=-L}^{L} |\phi_{Lm}|^2 ,
  \label{eq:block4}
\end{equation}
and this block resamples the $\phi$ prior spectrum every sweep and rebuilds the
Block 3 mass matrix to match, closing the joint posterior.

\subsubsection*{The prior on \texorpdfstring{$\clpp$}{ClPhiPhi}, and why it is
not a free choice}
\label{sec:clpp-prior}

Equation~\eqref{eq:block4} is the conditional implied by a \emph{flat} prior on
$\clpp$, which is improper. This is worth stating explicitly rather than
leaving implicit, because it propagates into the $\phi$ block in a way that is
easy to miss. Integrating $\clpp$ out of the joint gives a marginal on the
field, $p(\phi) \propto S_L^{-(L-1/2)}$; against the radial measure for the
$2L+1$ real components of $\phi_{Lm}$, $r^{2L}\,\mathrm{d}r$ with $S_L = r^2$,
this is
\begin{equation}
  p(r) \;\propto\; r^{\,1} ,
\end{equation}
i.e.\ flat in $S_L$ --- improper, and \emph{increasing} with amplitude. The
$\phi$ amplitude is therefore constrained by the lensing likelihood alone, with
no prior restoring force. Equivalently, substituting the conditional mean of
Eq.~\eqref{eq:block4} into the Gaussian $\phi$ prior term gives
$\tfrac12 \sum_L S_L/\clpp = \sum_L (L - 3/2)$, independent of $\phi$: the
prior is exactly scale-free along the amplitude ray.

In practice this is benign whenever the lensing signal-to-noise is sufficient
to pin the amplitude, which it is in the configurations reported here. It is
not benign as a matter of principle, and it has two concrete consequences that
we report rather than gloss: it removes the only safeguard against an
amplitude runaway if the temperature block is poorly initialised, and it means
a rank test of $\phi$ against a truth drawn from a \emph{proper} Gaussian prior
is not a calibration test of this target (Sec.~\ref{sec:sbc}).

We therefore also provide the conjugate proper alternative, a weakly
informative $\mathrm{InvGamma}(\nu/2,\, \nu C_L^{\phi\phi,\,\rm fid}/2)$ prior centred on
a fiducial spectrum, under which
\begin{equation}
  \clpp \mid \phi \;\sim\; \mathrm{InvGamma}\!\left(L - \tfrac12 + \tfrac{\nu}{2},\;
  \big(S_L + \nu\,C_L^{\phi\phi,\,\rm fid}\big)/2\right).
  \label{eq:block4-proper}
\end{equation}
Here $\nu$ reads as an effective number of prior pseudo-modes, on the same
footing as the $2L+1$ real modes the data supplies at multipole $L$, so the
prior regularises the lowest multipoles --- where the improper prior does its
damage --- and is negligible at high $L$. The marginal tail becomes
$p(r) \propto r^{\,1-\nu}$, so the target is proper only for $\nu > 2$; we use
$\nu = \placeholder{6}$, and reject $\nu \le 2$ at the API level rather than
allowing a choice that appears to fix the problem while leaving it in place.

% --- Forward model ---
\subsection{Differentiable lensing operator}
\label{sec:lensing}

The lensing forward model maps $(\alm, \phi) \to T_{\rm lensed}$ via bilinear
interpolation on the sphere, validated against the \texttt{healpy} deflection
reference and the dense Jacobian of \citet{Reinecke:2023} to machine precision.
Gradients with respect to both $\alm$ and $\phi$ are available via
\texttt{tf.custom\_gradient} / \texttt{tf.py\_function} wrapping of the ducc0
SHT; this is the same approach used in \citet{Flinch} and is treated here
as a standard building block.

% --- What the bias is NOT ---
\subsection{What the lensing-potential deficit is not}
\label{sec:deficit-not}

% ROADMAP.md: "Ruled out by construction — say so in three sentences in the
% paper, don't spend compute."
Several possible sources of systematic bias in the recovered $\phi$ posterior
are ruled out by the construction of this sampler and need not be tested
empirically.
(i)~\textit{N0/N1 noise bias}: a Gibbs sampler has no noise-bias subtraction
step; any bias in the recovered $\phi$ amplitude cannot originate from the
multiplicative noise floor of a quadratic estimator.
(ii)~\textit{Non-Gaussian deflections}: our validation simulations draw $\phi$
from its Gaussian prior exactly; non-Gaussianity of the lensing field is not a
factor.
(iii)~\textit{Foreground contamination}: the validation uses foreground-free
CMB-only simulations.
(iv)~\textit{Lensing-operator accuracy}: the matrix-free SHT forward model is
validated against the dense reference and independent \texttt{healpy} outputs to
numerical precision, following the standard set by \citet{Reinecke:2023}.
The residual $\phi$-amplitude deficit observed at finite chain length
(\S\ref{sec:results}) is consistent with an under-mixed / Wiener-suppressed
chain start \citep{Millea:2020}; it decreases monotonically with signal-to-noise
ratio ($\rho = -0.78$ across 27 multipole bins), which is the expected signature
of a sampler that has not yet equilibrated from a Wiener-filtered starting point.


% ============================================================
\section{Validation}
\label{sec:results}
% STATUS: [PLACEHOLDER — content depends on the equilibrated chains]

\subsection{Simulation-based calibration (SBC)}
\label{sec:sbc}

We run 12 independent chains at $\lmax = \placeholder{64}$, each with its own
$(\alm^{\rm true}, \phi^{\rm true}, \mathrm{noise})$ triple drawn from the
fiducial spectra and warm-started away from the truth, and rank the truth's
binned power within each posterior.

\subsubsection*{The rank statistic requires its own null}

The spectrum-level statistic must be calibrated before any departure from
uniformity is read as bias, and this is not a technicality. Ranking the truth's
\emph{realized} power $S_L/(2L+1)$ against posterior draws of $\clpp$ compares
it against the \emph{mode} of Eq.~\eqref{eq:block4}, since
$\beta/(\alpha+1) = S_L/(2L+1)$. The inverse-Gamma is right-skewed, so
$P(\text{draw} < \text{mode}) < 1/2$ \emph{for an exact sampler}, and averaging
over $\ell$ within a bin shrinks the spread while preserving the offset, driving
the mean rank toward zero. A naive uniformity test therefore rejects a correct
sampler with arbitrarily high confidence at large bin width.

We accordingly compare each observed rank against a null simulated from the
conditional itself. For $\cl^{TT}$ the observed bin means
$\placeholder{0.094, 0.104, 0.094, 0.375}$ agree with the simulated null
$\placeholder{0.095, 0.095, 0.115, 0.345}$ in every bin. For $\clpp$ the null
must additionally retain the chain's own sweep-to-sweep scatter in $\phi$,
which contributes a common-mode variance across $\ell$ that does not average
down within a bin; with that included, the observed
$\placeholder{0.281, 0.323, 0.323, 0.229}$ again lies within the null band
throughout. We report these rows as interval coverage, not calibration, since
the implied $\cl$ prior is flat and improper and no strict rank exists.

\subsubsection*{Block-level exactness}

Independently of the rank test, we verify Block 4 directly on the production
chains: using each sweep's $(\phi, \clpp)$ pair, the probability integral
transform $u = F_{\mathrm{InvGamma}(L-1/2,\,S_L(\phi)/2)}(\clpp)$ is uniform
(\placeholder{$\mathrm{KS}\ p = 0.53$}, \placeholder{$N = 17856$}), while a
deliberately lag-shifted control is rejected --- confirming the test has power
rather than merely failing to reject.

\subsubsection*{Field-level calibration}

The field ranks are the statistic this design supports directly. The $\alm$
ranks are consistent with uniformity. The $\phi$ ranks are not
(\placeholder{$\bar u = 0.367$}), and we attribute this to the improper prior
of Sec.~\ref{sec:clpp-prior} rather than to the sampler: with the flat prior
the target's $\phi$ marginal is flat in $S_L$, while $\phi^{\rm true}$ was drawn
from a proper Gaussian, so the two distributions differ by construction and the
rank has no reason to be uniform. We resolve this in two ways, and report both:
by fixing $\clpp$ at the fiducial spectrum, under which the $\phi$ prior is
proper and identical to the generative process; and by adopting the proper
conjugate prior of Eq.~\eqref{eq:block4-proper} and drawing the truth from that
same joint prior, which restores a strict SBC test while retaining the $\clpp$
block that produces the joint posterior of Sec.~\ref{sec:joint-posterior}.
\todo{Insert final rank figures and $p$-values from both runs.}

\subsection{Joint $(\cl^{TT}, \clpp)$ posterior}
\label{sec:joint-posterior}

\placeholder{
Figure: 2D marginal $(\cl^{TT}, \clpp)$ posterior from the pooled coverage
chains, showing the cross-correlation between the two power spectra.
This is the headline result: no existing method produces this joint object.
Result: \todo{correlation amplitude and sign.}
}

\subsection{C$_\ell^{TT}$ bias vs lensing-blind analysis}
\label{sec:cl-bias}

\placeholder{
Figure: recovered $\cl^{TT}$ from the lensing-aware joint sampler vs the
Commander-style lensing-blind Gibbs (same simulation, same noise, no $\phi$
block). The lensing-blind chain fits the lensed-sky power spectrum as if it were
unlensed; the bias shows directly how much lensing power leaks into the
inferred $\cl^{TT}$.
Data already available: \texttt{results/analysis/lensing\_blind\_baseline\_lmax128.npz}
(run submitted 2026-08-10, job 11717687).
}

\subsection{Per-mode uncertainty propagation}
\label{sec:uncertainty}

\placeholder{
Figure: posterior width on $\phi_{\ell m}$ from the joint sampler vs the
posterior width one would assign under the marginal (QE) approximation.
Quantifies what joint sampling buys over a marginalised method.
}


% ============================================================
\section{Discussion and conclusion}
\label{sec:discussion}
% STATUS: [PARTIAL — framing is fixed; numbers are placeholders]

We have presented the first full-sky, curved-sky joint Gibbs sampler over
$(\alm, \cl, \phi, \clpp)$.
The sampler's defining output — the joint $(\cl^{TT}, \clpp)$ posterior with
propagated correlations — is not produced by any existing method:
Commander-class samplers \citep{Commander,Almanac,Flinch} and QE-based methods
\citep{Carron:2017} both produce either the map-and-spectra posterior
(lensing-blind) or a lensing constraint (marginalised), never both together.

\textbf{Honest scale statement.}
The demonstration is at $\lmax \approx \placeholder{128}$, where chains
equilibrate in $\sim \placeholder{XX}$h on a single GPU node.
Scaling to $\lmax \sim 1000$ (the Planck / SO lensing science scale) is the
principal avenue for follow-up work and is not addressed here.

\textbf{As a reference standard.}
A sample-exact posterior is the object that amortised or learned CMB lensing
posteriors \citep{Doeser:2024} ultimately get validated against.
We offer the joint $(\cl^{TT}, \clpp)$ posterior at demonstration scale as
that reference.

\textbf{Outlook.}
Immediate extensions include:
real Planck data (\placeholder{in prep});
$\Lambda$CDM parameter inference from the posterior $\cl^{TT}$ chains;
polarisation (TQU) and LiteBIRD delensing forecasts via the spin-2 SHT
extension of the current infrastructure.

% ============================================================
\section*{Acknowledgements}
\todo{Acknowledgements to be added.}

% ============================================================
% REFERENCES
% Use inline \bibitem entries until a proper .bib file is assembled.
\begin{thebibliography}{99}

\bibitem{Flinch}
Crespi, S., Bonici, M., Loureiro, A., et al.\ 2025,
\textit{Flinch: A Differentiable Framework for Field-Level Inference of Cosmological Parameters from Curved Sky Data},
arXiv:2510.26691

\bibitem{Almanac}
Sellentin, E., Loureiro, A., Whiteway, L., et al.\ 2023,
\textit{Almanac: Scalable all-sky HMC CMB analysis},
arXiv:2305.16134

\bibitem{Commander}
Eriksen, H.\ K., et al.\ 2004,
\textit{Power Spectrum Estimation from the WMAP Data Using the Master HEALPix Estimator},
ApJS, 155, 227

\bibitem{Carron:2017}
Carron, J.\ \& Lewis, A.\ 2017,
\textit{Maximum a posteriori CMB lensing reconstruction},
PRD, 96, 063510,
arXiv:1704.08230

\bibitem{MUSE}
Millea, M.\ \& Seljak, U.\ 2022,
\textit{Marginal Unbiased Score Expansion and application to CMB lensing},
PRD, 105, 103531,
arXiv:2112.09091

\bibitem{CMBLensing}
% CMBLensing.jl — cite the published paper (Table II numbers used for comparison)
Millea, M., et al.\ 2021,
arXiv:2012.00011

\bibitem{Robnik:2022}
Robnik, J., De Luca, G.\ B., Silverstein, E.\ \& Seljak, U.\ 2022,
\textit{Microcanonical Langevin Monte Carlo},
arXiv:2212.08549

\bibitem{Bayer:2023}
Bayer, A.\ E., Seljak, U.\ \& Modi, C.\ 2023,
\textit{Microcanonical Langevin Monte Carlo for field-level inference of CMB},
arXiv:2307.09504

\bibitem{Reinecke:2023}
Reinecke, M., Belkner, S.\ \& Carron, J.\ 2023,
\textit{Faster CMB lensing with ducc0},
arXiv:2304.10431

\bibitem{Millea:2020}
Millea, M., Anderes, E.\ \& Wandelt, B.\ D.\ 2020,
\textit{Bayesian delensing of the CMB},
PRD, 102, 123542,
arXiv:2002.00965

\bibitem{Doeser:2024}
Doeser, L.\ \& Jasche, J.\ 2024,
arXiv:2606.10023

\bibitem{Drouin:2023}
% Score-based / diffusion generative lensing — representative reference
\todo{Fill in the correct diffusion-lensing citation from literature.md}

\end{thebibliography}

\end{document}
